\documentclass{applemlr}

\usepackage{amsmath, amssymb}
\usepackage{graphicx}
\usepackage{booktabs}
\usepackage{tabularx}
\usepackage{array}
\usepackage{enumitem}
\usepackage{listings}
\usepackage{caption}
\usepackage{float}

\setcounter{tocdepth}{2}
\setlist[itemize]{leftmargin=*, itemsep=0.4em}
\setlist[enumerate]{leftmargin=*, itemsep=0.4em}
\newcolumntype{Y}{>{\centering\arraybackslash}X}

\definecolor{listingbg}{gray}{0.95}
\definecolor{listingborder}{HTML}{D0D3D4}
\definecolor{listingcomment}{HTML}{5E8050}
\definecolor{listingkeyword}{HTML}{007AFF}

\lstdefinestyle{aiproject}{
  backgroundcolor=\color{listingbg},
  frame=single,
  framerule=0.4pt,
  rulecolor=\color{listingborder},
  basicstyle=\ttfamily\footnotesize,
  keywordstyle=\color{listingkeyword}\bfseries,
  commentstyle=\color{listingcomment},
  numbers=left,
  numberstyle=\scriptsize\color{listingkeyword},
  stepnumber=1,
  numbersep=6pt,
  breaklines=true,
  columns=fullflexible,
  captionpos=b,
  keepspaces=true
}


\usepackage{xcolor}       % 用于颜色定义
\usepackage[most]{tcolorbox} % 引入 tcolorbox 宏包

\newtcbox{\code}{
  nobeforeafter, % 确保它是一个行内命令
  colback=black!10, % 背景颜色 (10%的黑色，即浅灰色)
  colframe=black!60, % 边框颜色 (60%的黑色)
  boxrule=0.5pt,    % 边框粗细
  arc=2pt,          % 圆角大小
  boxsep=0pt,       % 内部文字与边框的间距
  left=3pt, right=3pt, top=2pt, bottom=2pt, % 分别设置四个方向的间距
  % raise=1.5pt, % 稍微向上调整盒子位置，让它和文本对齐更好看
  fontupper=\ttfamily % 盒子内的文字使用打字机字体
}




\title{课程大作业中文报告模板}

\author{张三}
\author{李四}
\author{王五}
\affiliation{上海交通大学自动化与感知学院}

\metadata[课程信息]{人工智能算法实践}
\metadata[提交日期]{2025 年 12 月 31 日}

\abstract{
本文提供一个简洁的人工智能课程大作业中文报告模板，帮助同学们快速上手撰写项目总结。模板涵盖文档结构建议、排版规范、图表与实验记录示例以及提交前检查清单，读者可直接在此基础上替换内容并生成高质量报告。
}

\begin{document}

\maketitle

\tableofcontents
\newpage

\section{快速开始}
\label{sec:quickstart}

\begin{enumerate}
  \item \textbf{复制模板：} 克隆仓库或复制本目录，在 \code{main.tex} 中替换作者信息、摘要和正文。
  \item \textbf{编译方式：} 推荐运行 \code{xelatex main.tex}；若添加参考文献，请依次执行 \code{bibtex main} 与额外两次 \code{xelatex main.tex} 以更新引用。
  \item \textbf{字体要求：} 模板默认启用 SF Pro 与自动检测的中文无衬线字体；若本地缺失，请安装 \code{Noto Sans CJK SC} 或 \code{Source Han Sans SC}。
  \item \textbf{参考文献：} 所有条目集中维护在项目根目录的 \code{ref.bib} 中，引用方式参见附录~\ref{sec:bib}。
  \item \textbf{代码高亮：} 使用 \code{listings} 宏包，示例参见\autoref{lst:code}；如需 \code{minted} 请自行启用并添加 \code{-shell-escape}。
  \item \textbf{提交产物：} 建议提交 \code{main.pdf} 以及必要的源文件、数据和运行脚本。
\end{enumerate}

\section{建议的报告结构}
\label{sec:structure}

\subsection{封面与摘要}
\begin{itemize}
  \item \textbf{标题：} 明确描述项目主题，例如 ``基于 Transformer 的中文情感分析''。
  \item \textbf{作者与学号：} 建议使用 \code{\string\author} 和 \code{\string\affiliation} 列出全部成员及分工说明。
  \item \textbf{摘要：} 150--200 字，概述任务背景、方法要点、核心结果与结论。
\end{itemize}

\subsection{正文主体}
可按以下层次组织内容（可增删）：
\begin{enumerate}
  \item \textbf{引言：} 说明问题背景、研究动机与项目目标。
  \item \textbf{相关工作：} 简述参考文献或基线方法，突出差异。
  \item \textbf{方法设计：} 介绍模型架构、算法流程和实现细节，可适量插入公式。
  \item \textbf{实验设置：} 给出数据集来源、预处理方式、评价指标和训练超参数。
  \item \textbf{结果与分析：} 使用表格、图像比较方法表现，讨论优势与不足。
  \item \textbf{总结与展望：} 回顾主要收获并提出后续改进方向。
\end{enumerate}

\section{图表与排版示例}
\label{sec:figures}

\subsection{图像示例}
\begin{figure}[ht]
  \centering
  \fbox{\parbox[c][0.4\linewidth][c]{0.8\linewidth}{\centering 在此放置实验结果示意图\\(可替换为真实图片文件)}}
  \caption{模型推理流程或实验可视化示例。上传正式报告时请替换为真实图像并更新说明。}
  \label{fig:placeholder}
\end{figure}

\subsection{表格示例}
\begin{table}[ht]
  \centering
  \caption{模型在验证集上的性能指标示例。}
  \label{tab:metrics}
  \begin{tabularx}{0.9\linewidth}{@{}l *{3}{Y}@{}}
    \toprule
    \textbf{模型} & \textbf{准确率} & \textbf{宏平均 F1} & \textbf{推理时间 (ms)} \\
    \midrule
    基线 SVM & 86.2 & 85.0 & 4.8 \\
    BERT 微调 & 92.7 & 92.1 & 16.4 \\
    我们的方法 & \textbf{94.5} & \textbf{94.0} & 18.9 \\
    \bottomrule
  \end{tabularx}
\end{table}

\subsection{代码示例}
\lstinputlisting[style=aiproject, language=Python, caption={训练脚本片段示例。}, label={lst:code}]{store/example.py}


\section{实验记录与可复现性}
\label{sec:repro}

\begin{itemize}
  \item \textbf{环境说明：} 记录硬件、操作系统、Python/框架版本，并给出 \code{requirements.txt} 或 \code{environment.yml}。
  \item \textbf{训练流程：} 说明数据划分、训练轮次、学习率、随机种子等关键参数。
  \item \textbf{可重复验证：} 提供运行命令示例，例如 \code{python train.py --config configs/base.yaml}，并展示日志或曲线截图。
\end{itemize}

\section{提交前检查清单}
\label{sec:checklist}

\begin{enumerate}
  \item 已完成全部小节内容更新，删除无关示例文字与占位图。
  \item 图表均有编号与清晰说明，坐标与单位准确无误。
  \item 公式、变量符号在全文保持一致，并适当给出解释。
  \item 引用格式统一，\code{ref.bib} 中条目完整无误。
  \item 运行 \code{xelatex main.tex} 至少两次，确保目录与引用无报错；若使用 \code{bibtex}，请执行额外编译流程。
  \item 清理辅助文件（可执行 \code{latexmk -c}），提交结果前手动检查 PDF 排版。
\end{enumerate}

\addcontentsline{toc}{section}{参考文献}
\bibliographystyle{ieeetr}
\bibliography{ref}

\appendix

\section{参考文献管理示例}
\label{sec:bib}

使用 \code{bibtex} 时可将文献条目放入 \code{ref.bib}，并在正文中通过 \code{\string\citep{lecun2015deep}} 引用。编译流程示意：
\begin{enumerate}
  \item 首次编译，立即暴露语法与缺字错误：\code{xelatex main.tex}
  \item 更新 \code{ref.bib} 后运行，生成最新参考文献列表：\code{bibtex main}
  \item 重新编译以刷新目录、交叉引用与页码：\code{xelatex main.tex} （运行两次）
  % \item 重新编译以刷新目录、交叉引用与页码：xelatex main.tex
\end{enumerate}



下面给出占位符文献条目（默认已写入 \code{ref.bib}）：
\begin{verbatim}
@article{lecun2015deep,
  title={Deep learning},
  author={LeCun, Yann and Bengio, Yoshua and Hinton, Geoffrey},
  journal={nature},
  volume={521},
  number={7553},
  pages={436--444},
  year={2015},
  publisher={Nature Publishing Group UK London}
}
\end{verbatim}

\section{常见问题解答}

\begin{description}[leftmargin=0.5cm,labelindent=0cm]
  \item[图像模糊] 使用矢量格式（PDF/SVG）或更高分辨率 PNG。
  \item[中文乱码] 确认以 UTF-8 编码保存 \code{main.tex}，并使用 XeLaTeX 编译。
\end{description}

\end{document}
